\chapter{Topologia da Reta Real}\label{ch:b1:topology}

Os limites remetem sempre ao mesmo vocabulário geométrico: pontos
``próximos de'' um \omterm{def:b1:logic:sets}{conjunto}, \omterm{def:b1:logic:sets}{conjuntos} ``sem vazamentos pela
\omterm{def:b1:topology:closure}{fronteira}'', \omterm{prop:b1:reals:intervals}{intervalos} dos quais as sequências não conseguem escapar.
Este capítulo fixa esse vocabulário --- \omterm{def:b1:topology:open}{conjuntos abertos} e \omterm{def:b1:topology:closed}{fechados},
\omterm{def:b1:topology:closure}{interior} e \omterm{def:b1:topology:closure}{fecho}, \omterm{def:b1:topology:dense}{densidade} --- na reta real, e demonstra a compacidade
dos segmentos na sua forma sequencial. As mesmas noções, em espaços
vetoriais normados, são matéria do segundo ano; em $\R$ elas estão ao
alcance e são imediatamente úteis para o \cref{ch:b1:continuity}.

\section{Conjuntos abertos, conjuntos fechados}

\begin{definition}[Vizinhança, conjunto aberto]\label{def:b1:topology:open}
Um \omterm{def:b1:logic:sets}{conjunto} $V \subseteq \R$ é uma \emph{vizinhança}\index{vizinhança}
de $x \in \R$ quando contém um \omterm{prop:b1:reals:intervals}{intervalo} $\intoo{x - r}{x + r}$ para
algum $r > 0$. Um \omterm{def:b1:logic:sets}{conjunto} $U \subseteq \R$ é \emph{aberto}\index{conjunto aberto}
quando é vizinhança de cada um dos seus pontos:
\[
\forall x \in U,\ \exists r > 0, \quad \intoo{x - r}{x + r}
\subseteq U .
\]
\end{definition}

\begin{example}\label{ex:b1:topology:open}
Os \omterm{prop:b1:reals:intervals}{intervalos} \omterm{def:b1:topology:open}{abertos} são \omterm{def:b1:topology:open}{abertos}: para $x \in \intoo{a}{b}$, tome $r = \min(x -
a,\, b - x) > 0$. As semirretas $\intoo{a}{+\infty}$ são \omterm{def:b1:topology:open}{abertas}; $\R$ e
$\emptyset$ são \omterm{def:b1:topology:open}{abertos} (o último por vacuidade). $\intcc{0}{1}$ \emph{não}
é \omterm{def:b1:topology:open}{aberto}: nenhum \omterm{prop:b1:reals:intervals}{intervalo} em torno de $0$ permanece dentro.
\end{example}

\begin{proposition}[Estabilidade dos abertos]\label{prop:b1:topology:openstable}
Toda união de \omterm{def:b1:topology:open}{abertos} é \omterm{def:b1:topology:open}{aberta}; uma interseção \emph{finita} de
\omterm{def:b1:topology:open}{abertos} é \omterm{def:b1:topology:open}{aberta}. Interseções infinitas podem falhar:
$\bigcap_{n \geq 1} \intoo{-\frac 1n}{\frac 1n} = \{0\}$, que não é \omterm{def:b1:topology:open}{aberto}.
\end{proposition}

\begin{proof}
União: se $x \in \bigcup_i U_i$, então $x \in U_{i_0}$ para algum $i_0$,
e o \omterm{prop:b1:reals:intervals}{intervalo} fornecido por $U_{i_0}$ cabe dentro da união. Interseção
finita: se $x \in U_1 \cap \dots \cap U_k$, tome $r = \min(r_1,
\dots, r_k) > 0$ dos raios fornecidos por cada $U_j$. Para o
contraexemplo: todo \omterm{prop:b1:reals:intervals}{intervalo} em torno de $0$ contém algum $\frac{1}{n}$
(Arquimedes) e, portanto, sai da interseção.
\end{proof}

\begin{example}[Certificando a abertura com raios explícitos]\label{ex:b1:topology:radii}
O \omterm{def:b1:logic:sets}{conjunto} $U = \{x \in \R : x^2 > 2\}$ é \omterm{def:b1:topology:open}{aberto}? Sim, e o certificado
pode ser escrito: $U = \intoo{-\infty}{-\sqrt2} \cup
\intoo{\sqrt2}{+\infty}$, uma união de duas semirretas \omterm{def:b1:topology:open}{abertas}, \omterm{def:b1:topology:open}{aberta}
pela \cref{prop:b1:topology:openstable}. Alternativamente, argumente
ponto a ponto: para $x \in U$ com $x > \sqrt 2$, tome $r = x -
\sqrt2 > 0$: todo $y \in \intoo{x - r}{x + r}$ satisfaz $y >
\sqrt 2$, logo $y^2 > 2$; simetricamente à esquerda. Os dois
estilos importam --- o estrutural (construir a partir de \omterm{def:b1:topology:open}{abertos}
conhecidos por uniões e interseções finitas) escala melhor, e o
estilo $\varepsilon$ funciona quando nenhuma estrutura é visível; e o
\cref{ch:b1:continuity} acrescentará um terceiro, o mais poderoso:
$U$ é a \omterm{def:b1:logic:map}{pré-imagem} do \omterm{def:b1:topology:open}{aberto} $\intoo{2}{+\infty}$ pela
\omterm{def:b1:logic:map}{aplicação} contínua $x \mapsto x^2$.
\end{example}

\begin{definition}[Conjunto fechado]\label{def:b1:topology:closed}
Um \omterm{def:b1:logic:sets}{conjunto} $F \subseteq \R$ é \emph{fechado}\index{conjunto fechado} quando o
seu complementar $\R \setminus F$ é \omterm{def:b1:topology:open}{aberto}. Por de Morgan e pela
\cref{prop:b1:topology:openstable}: toda interseção de fechados é
fechada, e uniões finitas de fechados são fechadas.
\end{definition}

\begin{theorem}[Caracterização sequencial dos fechados]\label{thm:b1:topology:seqclosed}
$F$ é \omterm{def:b1:topology:closed}{fechado} se, e somente se: para toda sequência $(u_n)$ de pontos
de $F$ que converge para algum $\ell \in \R$, o limite $\ell$ pertence a
$F$. (``\omterm{def:b1:topology:closed}{Fechado}'' $=$ ``estável por limites''.)
\end{theorem}

\begin{proof}
($\Rightarrow$) Seja $F$ \omterm{def:b1:topology:closed}{fechado}, $u_n \in F$, $u_n \to \ell$, e
suponha $\ell \notin F$. O complementar é \omterm{def:b1:topology:open}{aberto}: algum
$\intoo{\ell - r}{\ell + r}$ evita $F$. Mas a convergência põe $u_n$
nesse \omterm{prop:b1:reals:intervals}{intervalo} para $n$ grande: contradição com $u_n \in F$.

($\Leftarrow$) Suponha que $F$ não seja \omterm{def:b1:topology:closed}{fechado}: o complementar não é \omterm{def:b1:topology:open}{aberto},
de modo que algum $x \notin F$ não tem \omterm{prop:b1:reals:intervals}{intervalo} $\intoo{x - r}{x + r}$ dentro
do complementar; tomando $r = \frac{1}{n+1}$, escolha $u_n \in F$ com
$\abs{u_n - x} < \frac{1}{n+1}$. Então $u_n \in F$, $u_n \to x \notin
F$: a propriedade sequencial falha.
\end{proof}

\begin{example}[O teste sequencial, nos dois sentidos]\label{ex:b1:topology:seqtest}
\emph{\omterm{def:b1:topology:closed}{Fechado}:} $F = \Z \cup \bigl\{n + \frac1n : n \geq 2\bigr\}$.
Seja $u_k \in F$ com $u_k \to \ell$. A janela $\intcc{\ell -
1}{\ell + 1}$ contém apenas finitos pontos de $F$ (finitos inteiros, finitos
$n + \frac1n$), e, a partir de certo posto, todos os $u_k$ estão nela:
a sequência assume então finitos valores e, convergindo, é
constante a partir de certo ponto (como no
\cref{exo:b1:topology:3}): $\ell \in F$. \omterm{def:b1:topology:closed}{Fechado} --- embora
$F$ contenha pares de pontos a distância $\frac1n$, arbitrariamente
próximos.

\emph{Não \omterm{def:b1:topology:closed}{fechado}:} $G = \bigl\{\frac1m + \frac1n : m, n \in
\N^*\bigr\}$. A sequência $\frac1n + \frac1n \in G$ tende a
$0$, e $0 \notin G$ (soma de dois termos positivos): o
teste sequencial falha, e $G$ não é \omterm{def:b1:topology:closed}{fechado}. Curiosamente, cada
$\frac1m$ \emph{pertence} a $\overline G \cap G$: com efeito,
$\frac1m = \frac{1}{m+1} + \frac{1}{m(m+1)} \in G$. A ideia de
fechamento: para demonstrar que um \omterm{def:b1:logic:sets}{conjunto} é \omterm{def:b1:topology:closed}{fechado}, controle
\emph{todas} as sequências convergentes de uma vez (em geral por
finitude local ou por um argumento de fórmula \omterm{def:b1:topology:closed}{fechada}); para refutá-lo,
basta uma sequência de fuga bem escolhida --- a assimetria torna a
direção negativa a fácil, e os contraexemplos deste capítulo têm todos
essa forma de uma linha.
\end{example}

\begin{example}\label{ex:b1:topology:closed}
Os segmentos $\intcc{a}{b}$, as semirretas $\intco{a}{+\infty}$, os \omterm{def:b1:counting:card}{conjuntos
finitos} e $\Z$ (uma sequência convergente de inteiros é constante a
partir de certo ponto) são \omterm{def:b1:topology:closed}{fechados}. $\intoc{0}{1}$ não é \omterm{def:b1:topology:open}{aberto} (falha em
$1$) nem \omterm{def:b1:topology:closed}{fechado} ($\frac 1n \to 0 \notin$ do \omterm{def:b1:logic:sets}{conjunto}): a maioria dos \omterm{def:b1:logic:sets}{conjuntos} não
é nem uma coisa nem outra. $\R$ e $\emptyset$ são ao mesmo tempo \omterm{def:b1:topology:open}{abertos}
e \omterm{def:b1:topology:closed}{fechados} --- e são os únicos subconjuntos de $\R$ assim
(\cref{exo:b1:topology:9}).
\end{example}

\begin{example}[Um aberto montado a partir de infinitas peças]\label{ex:b1:topology:rminusz}
$\R \setminus \Z = \bigcup_{n \in \Z} \intoo{n}{n+1}$: uma
união infinita de \omterm{prop:b1:reals:intervals}{intervalos} \omterm{def:b1:topology:open}{abertos}, \omterm{def:b1:topology:open}{aberta} pela
\cref{prop:b1:topology:openstable} --- de modo que $\Z$ é \omterm{def:b1:topology:closed}{fechado} sem
necessidade de argumento sequencial algum. Note a divisão de trabalho
nas regras de estabilidade: as \emph{uniões} de \omterm{def:b1:topology:open}{abertos} podem ser
arbitrárias (cada ponto só precisa do seu próprio certificado,
fornecido pelo único \omterm{def:b1:logic:sets}{conjunto} que o contém), ao passo que as
\emph{interseções} devem permanecer finitas (os certificados têm de ser
interceptados, e infinitos raios podem encolher a nada). O
\Cref{exo:b1:topology:10} mostrará que este exemplo é a forma geral: todo
subconjunto \omterm{def:b1:topology:open}{aberto} de $\R$ é uma união enumerável e disjunta de
\omterm{prop:b1:reals:intervals}{intervalos} \omterm{def:b1:topology:open}{abertos}.
\end{example}

\section{Interior, fecho, densidade}

\begin{definition}[Interior, fecho, fronteira]\label{def:b1:topology:closure}
Seja $A \subseteq \R$.
\begin{itemize}
  \item Um ponto $x$ é \emph{interior} a $A$ quando $A$ é uma
        \omterm{def:b1:topology:open}{vizinhança} de $x$; o \emph{interior}
        $\mathring{A}$\index{interior} é o \omterm{def:b1:logic:sets}{conjunto} dos pontos interiores.
  \item Um ponto $x$ é \emph{aderente} a $A$ quando toda \omterm{def:b1:topology:open}{vizinhança}
        de $x$ encontra $A$; o \emph{fecho}
        $\overline{A}$\index{fecho} é o \omterm{def:b1:logic:sets}{conjunto} dos pontos aderentes.
  \item A \emph{fronteira} é $\partial A = \overline A \setminus
        \mathring A$.\index{fronteira}
\end{itemize}
Então $\mathring A \subseteq A \subseteq \overline A$.
\end{definition}

\begin{proposition}[Propriedades principais]\label{prop:b1:topology:closureprops}
\begin{enumerate}
  \item $\mathring A$ é o maior \omterm{def:b1:topology:open}{aberto} contido em $A$; $A$
        é \omterm{def:b1:topology:open}{aberto} se, e somente se, $A = \mathring A$.
  \item $\overline A$ é o menor \omterm{def:b1:topology:closed}{fechado} que contém $A$; $A$
        é \omterm{def:b1:topology:closed}{fechado} se, e somente se, $A = \overline A$.
  \item (Caracterização sequencial da aderência) $x \in \overline
        A$ se, e somente se, $x$ é limite de uma sequência de pontos de
        $A$.
  \item A complementação troca as noções: $\R \setminus
        \overline A = \bigl(\R \setminus A\bigr)^{\!\circ}$.
\end{enumerate}
\end{proposition}

\begin{proof}
(4) $x \notin \overline A$ $\iff$ alguma \omterm{def:b1:topology:open}{vizinhança} de $x$ evita $A$
$\iff$ algum \omterm{prop:b1:reals:intervals}{intervalo} em torno de $x$ está em $\R \setminus A$ $\iff$ $x$
é \omterm{def:b1:topology:closure}{interior} a $\R \setminus A$.

(1) $\mathring A$ é \omterm{def:b1:topology:open}{aberto}: se $x \in \mathring A$, algum
$\intoo{x-r}{x+r} \subseteq A$; todo ponto $y$ desse \omterm{prop:b1:reals:intervals}{intervalo} tem um \omterm{prop:b1:reals:intervals}{intervalo}
menor em torno de si dentro dele, logo dentro de $A$: o \omterm{prop:b1:reals:intervals}{intervalo}
inteiro está em $\mathring A$. Todo \omterm{def:b1:topology:open}{aberto} $U \subseteq A$ é formado por
pontos \omterm{def:b1:topology:closure}{interiores} de $A$, de modo que $U \subseteq \mathring A$: é o maior. A
caracterização da abertura decorre daí.

(2) Em detalhe. Por (4), $\R \setminus \overline A$ é o \omterm{def:b1:topology:closure}{interior}
de $\R \setminus A$, um \omterm{def:b1:topology:open}{aberto} por (1): logo, $\overline A$ é
\omterm{def:b1:topology:closed}{fechado}, e contém $A$. Minimalidade: seja $F \supseteq A$ \omterm{def:b1:topology:closed}{fechado}.
Então $\R \setminus F$ é \omterm{def:b1:topology:open}{aberto} e está contido em $\R
\setminus A$, de modo que, pela maximalidade em (1),
\[
\R \setminus F \subseteq \bigl(\R \setminus A\bigr)^{\!\circ}
= \R \setminus \overline A ,
\]
e, tomando complementares de novo: $\overline A \subseteq F$. Assim,
$\overline A$ é o menor \omterm{def:b1:topology:closed}{fechado} que contém $A$. Caracterização:
se $A = \overline A$, então $A$ é \omterm{def:b1:topology:closed}{fechado} (acabamos de ver); se $A$ é
\omterm{def:b1:topology:closed}{fechado}, ele próprio é um \omterm{def:b1:topology:closed}{fechado} que contém $A$, de modo que a
minimalidade força $\overline A \subseteq A$, e daí a igualdade.

(3) Se $u_n \in A$, $u_n \to x$: toda \omterm{def:b1:topology:open}{vizinhança} de $x$ contém
algum $u_n \in A$, de sorte que $x \in \overline A$. Reciprocamente, se $x \in
\overline A$: cada \omterm{prop:b1:reals:intervals}{intervalo} $\intoo{x - \frac{1}{n+1}}{x +
\frac{1}{n+1}}$ encontra $A$ em algum $u_n$, e $u_n \to x$.
\end{proof}

\begin{example}\label{ex:b1:topology:closureexamples}
$\overline{\intoo{0}{1}} = \intcc{0}{1}$;
$\mathring{\intcc{0}{1}} = \intoo{0}{1}$; $\partial\intoo{0}{1} =
\{0, 1\}$. Para $A = \{\frac 1n : n \in \N^*\}$: $\overline A = A
\cup \{0\}$, $\mathring A = \emptyset$, $\partial A = A \cup \{0\}$.
Para $\Q$: pela \omterm{def:b1:topology:dense}{densidade} (\cref{thm:b1:reals:density}), todo real é
aderente a $\Q$, de modo que $\overline{\Q} = \R$, ao passo que $\mathring{\Q} =
\emptyset$ (todo \omterm{prop:b1:reals:intervals}{intervalo} contém irracionais): a \omterm{def:b1:topology:closure}{fronteira} de
$\Q$ é todo o $\R$.
\end{example}

\begin{example}[Uma anatomia completa]\label{ex:b1:topology:anatomy}
Seja $A = \intoc{0}{1} \,\cup\, \bigl(\Q \cap \intoo{2}{3}\bigr)
\,\cup\, \{4\}$. Calculamos os três \omterm{def:b1:logic:sets}{conjuntos} da
\cref{def:b1:topology:closure}, peça por peça.

\emph{\omterm{def:b1:topology:closure}{Interior}.} Um ponto de $\intoo{0}{1}$ tem um \omterm{prop:b1:reals:intervals}{intervalo} inteiro
dentro de $A$: é \omterm{def:b1:topology:closure}{interior}. O ponto $1$: todo \omterm{prop:b1:reals:intervals}{intervalo} em torno dele
vaza à direita de $1$, onde $A$ nada tem até $2$: não é
\omterm{def:b1:topology:closure}{interior}. Nenhum ponto de $\Q \cap \intoo{2}{3}$ é \omterm{def:b1:topology:closure}{interior} (todo
\omterm{prop:b1:reals:intervals}{intervalo} contém irracionais, \cref{thm:b1:reals:density});
nem o ponto isolado $4$. Logo, $\mathring A = \intoo{0}{1}$.

\emph{\omterm{def:b1:topology:closure}{Fecho}.} Limites de pontos de $A$: todo $\intcc{0}{1}$
($0 = \lim \frac1n$, com $\frac 1n \in A$); todo
$\intcc{2}{3}$ (todo real ali é limite de racionais do
\omterm{prop:b1:reals:intervals}{intervalo}, \omterm{def:b1:topology:dense}{densidade} de novo); e $4$. Nada mais: um ponto
fora de $\intcc{0}{1} \cup \intcc{2}{3} \cup \{4\}$ tem distância positiva a
esse \omterm{def:b1:topology:closed}{conjunto fechado}. Logo, $\overline A = \intcc{0}{1} \cup
\intcc{2}{3} \cup \{4\}$.

\emph{\omterm{def:b1:topology:closure}{Fronteira}.} $\partial A = \overline A \setminus \mathring A
= \{0, 1\} \cup \intcc{2}{3} \cup \{4\}$.

A ideia de fechamento: as três operações agem \emph{localmente} ---
cada peça de $A$ contribui conforme a sua própria natureza (um
\omterm{prop:b1:reals:intervals}{intervalo} maciço guarda o seu \omterm{def:b1:topology:closure}{interior}, uma peça \omterm{def:b1:topology:dense}{densa} mas porosa se
transforma inteiramente em \omterm{def:b1:topology:closure}{fronteira}, e um ponto isolado é pura
\omterm{def:b1:topology:closure}{fronteira}) --- e um desenho de duas linhas de $A$ prevê todas as
respostas antes de escrita qualquer demonstração.
\end{example}

\begin{remark}[Armadilhas frequentes no raciocínio conjuntista]
(i) \emph{``Não \omterm{def:b1:topology:open}{aberto}'' não significa ``\omterm{def:b1:topology:closed}{fechado}''}: a maioria dos
\omterm{def:b1:logic:sets}{conjuntos} não é nem uma coisa nem outra ($\intoc{0}{1}$), e dois \omterm{def:b1:logic:sets}{conjuntos}
são as duas coisas ($\emptyset$, $\R$) --- \omterm{def:b1:topology:open}{aberto} e \omterm{def:b1:topology:closed}{fechado} não são
opostos, e sim duais por complementação. (ii) \emph{\omterm{def:b1:topology:closure}{Interior} e \omterm{def:b1:topology:closure}{fecho}
não comutam}: para $A = \Q$,
\[
\overline{\mathring A} = \overline\emptyset = \emptyset
\qquad\text{ao passo que}\qquad
\bigl(\,\overline A\,\bigr)^{\!\circ} = \mathring \R = \R :
\]
os dois operadores iterados diferem o máximo que \omterm{def:b1:logic:sets}{conjuntos} podem
diferir. (iii) \emph{Uniões infinitas de \omterm{def:b1:topology:closed}{fechados} podem falhar em ser
\omterm{def:b1:topology:closed}{fechadas}}: $\bigcup_{n\geq1} \intcc{\frac1n}{1} = \intoc{0}{1}$ --- o
espelho do contraexemplo da interseção da
\cref{prop:b1:topology:openstable}. (iv) \emph{\omterm{def:b1:topology:dense}{Denso} não significa
grande}: $\Q$ é \omterm{def:b1:topology:dense}{denso}, enumerável, de \omterm{def:b1:topology:closure}{interior} vazio, e o seu
complementar também é \omterm{def:b1:topology:dense}{denso}; a \omterm{def:b1:topology:dense}{densidade} diz ``arbitrariamente perto de
tudo'', e não ``quase tudo'' --- o \omterm{pb:b1:topology:1}{conjunto de Cantor} do problema de
fim de semana (\cref{pb:b1:topology:1}) faz o ponto oposto: um
\omterm{def:b1:logic:sets}{conjunto} topologicamente pequeno que é imenso em \omterm{def:b1:counting:card}{cardinalidade}.
\end{remark}

\begin{example}[Um fecho calculado exatamente]\label{ex:b1:topology:closurecomputed}
Seja $G = \bigl\{\frac1m + \frac1n : m, n \in \N^*\bigr\}$ (do
\cref{ex:b1:topology:seqtest}). Afirmação:
\[
\overline G = G \,\cup\, \Bigl\{\frac1m : m \in \N^*\Bigr\}
\,\cup\, \{0\} .
\]
($\supseteq$) $\frac1m = \lim_n \bigl(\frac1m + \frac1n\bigr)$
e $0 = \lim_n \frac2n$: aderentes, pela caracterização
sequencial. ($\subseteq$) Seja $x = \lim_k \bigl(
\frac{1}{m_k} + \frac{1}{n_k}\bigr)$; ordene cada par de modo que
$m_k \leq n_k$. Se $(m_k)$ é ilimitada, uma \omterm{def:b1:seq:subsequence}{subsequência} tem $m_k
\to \infty$, donde $n_k \to \infty$ também, e $x = 0$. Caso contrário,
$(m_k)$ assume finitos valores, um deles, digamos $m$,
infinitas vezes; ao longo dessa \omterm{def:b1:seq:subsequence}{subsequência}, $\frac{1}{n_k} \to x -
\frac1m$: se $(n_k)$ é limitada, ela assume algum valor $n$
infinitas vezes e $x = \frac1m + \frac1n \in G$; se não, $x =
\frac1m$. Todos os casos caem no \omterm{def:b1:logic:sets}{conjunto} anunciado. A ideia de
fechamento: calcular um \omterm{def:b1:topology:closure}{fecho} é uma análise de casos ao estilo da
compacidade sobre os índices --- índice limitado significa finitos
valores (casa dos pombos), índice ilimitado significa um limite que
escapa --- e a resposta exibe a típica estrutura em duas camadas dos
pontos-limite: o \omterm{def:b1:logic:sets}{conjunto}, os seus limites de primeira geração e o
limite deles, $0$.
\end{example}

\begin{definition}[Densidade, forma topológica]\label{def:b1:topology:dense}
$A$ é \emph{denso}\index{densidade} em $\R$ quando $\overline A = \R$
--- equivalentemente, todo \omterm{prop:b1:reals:intervals}{intervalo} \omterm{def:b1:topology:open}{aberto} não vazio encontra $A$;
equivalentemente (pela \cref{prop:b1:topology:closureprops}~(3)), todo
real é limite de elementos de $A$. Exemplos: $\Q$, $\R \setminus
\Q$,
os diádicos (\cref{exo:b1:reals:8}), os \omterm{def:b1:structures:subgroup}{subgrupos} densos
(\cref{exo:b1:reals:9}).
\end{definition}

\begin{example}[A densidade é relativa]\label{ex:b1:topology:relativedensity}
``\omterm{def:b1:topology:dense}{Denso}'', como definido aqui, significa \omterm{def:b1:topology:dense}{denso} \emph{em $\R$}; um
\omterm{def:b1:logic:sets}{conjunto} pode, em vez disso, ser \omterm{def:b1:topology:dense}{denso} apenas numa parte da reta. Os
diádicos de $\intcc{0}{1}$, isto é, $D \cap \intcc{0}{1}$
(\cref{exo:b1:reals:8}), encontram todo \omterm{prop:b1:reals:intervals}{intervalo} \omterm{def:b1:topology:open}{aberto} contido em
$\intcc{0}{1}$, mas evidentemente não tocam $\intoo{2}{3}$: eles são
\omterm{def:b1:topology:dense}{densos} \emph{em} $\intcc{0}{1}$, o que significa $\overline{D \cap
\intcc{0}{1}} = \intcc{0}{1}$. A expressão geral ``$A$ é \omterm{def:b1:topology:dense}{denso}
em $B$'' abrevia $B \subseteq \overline A$ --- nomeie sempre o \omterm{def:b1:logic:sets}{conjunto}
ambiente, pois os extremos do problema de fim de semana são \omterm{def:b1:topology:dense}{densos}
\emph{no \omterm{pb:b1:topology:1}{conjunto de Cantor}}, sendo nunca \omterm{def:b1:topology:dense}{densos} \emph{em
$\R$}: o mesmo \omterm{def:b1:logic:sets}{conjunto}, duas descrições verdadeiras e de som oposto.
\end{example}

\begin{example}[Manejando a densidade]\label{ex:b1:topology:densitytricks}
Três movimentos rápidos que recorrem constantemente. \emph{Ampliar}: se
$A$ é \omterm{def:b1:topology:dense}{denso} e $A \subseteq B$, então $B$ é \omterm{def:b1:topology:dense}{denso} (todo
\omterm{prop:b1:reals:intervals}{intervalo} já encontra $A$). \emph{Transportar}: se $A$ é
\omterm{def:b1:topology:dense}{denso}, então $\lambda A + \mu$ também o é, para $\lambda \neq 0$ --- um
\omterm{prop:b1:reals:intervals}{intervalo} $I$ encontra $\lambda A + \mu$ se, e somente se, o \omterm{prop:b1:reals:intervals}{intervalo}
$\frac{I - \mu}{\lambda}$ encontra $A$; assim, os múltiplos ímpares de
$10^{-9}$, digamos, são \omterm{def:b1:topology:dense}{densos}. \emph{Interceptar falha}: dois \omterm{def:b1:logic:sets}{conjuntos}
\omterm{def:b1:topology:dense}{densos} podem não se tocar ($\Q$ e $\R \setminus \Q$):
a \omterm{def:b1:topology:dense}{densidade} sobrevive a uniões e a aplicações afins, jamais a
interseções.
\end{example}

\section{Compacidade dos segmentos}

\begin{theorem}[Os segmentos são sequencialmente compactos]\label{thm:b1:topology:compact}
Seja $a \leq b$. Toda sequência de pontos de $\intcc{a}{b}$ tem uma
\omterm{def:b1:seq:subsequence}{subsequência} que converge para um ponto \emph{de}
$\intcc{a}{b}$.\index{compacidade}

Mais geralmente, os subconjuntos de $\R$ com essa propriedade (toda
sequência tem uma \omterm{def:b1:seq:subsequence}{subsequência} que converge no \omterm{def:b1:logic:sets}{conjunto}) são exatamente
os \omterm{def:b1:logic:sets}{conjuntos} \emph{\omterm{def:b1:topology:closed}{fechados} e limitados}.
\end{theorem}

\begin{proof}
Uma sequência em $\intcc{a}{b}$ é limitada, de modo que
Bolzano--Weierstrass (\cref{thm:b1:seq:bw}) extrai uma \omterm{def:b1:seq:subsequence}{subsequência}
convergente; o seu limite permanece em $\intcc{a}{b}$, pois os segmentos
são \omterm{def:b1:topology:closed}{fechados} (\cref{thm:b1:topology:seqclosed}).

Caso geral. \emph{(\omterm{def:b1:topology:closed}{Fechado} e limitado $\Rightarrow$ compacto)}: sejam
$F$ \omterm{def:b1:topology:closed}{fechado} e limitado e $(u_n)$ uma sequência em $F$.
A limitação de $F$ limita a sequência, de modo que
Bolzano--Weierstrass extrai $u_{\varphi(n)} \to \ell$; e $\ell \in F$, pois $F$
é \omterm{def:b1:topology:closed}{fechado} e a \omterm{def:b1:seq:subsequence}{subsequência} é uma sequência convergente de pontos
de $F$ (\cref{thm:b1:topology:seqclosed}): as duas hipóteses são
consumidas uma a uma, a limitação para a existência do limite e o
fechamento para a sua pertinência. \emph{(Compacto $\Rightarrow$ \omterm{def:b1:topology:closed}{fechado} e
limitado)}: se $F$ é ilimitado, tome $u_n \in F$ com $\abs{u_n} \geq
n$; toda \omterm{def:b1:seq:subsequence}{subsequência} é ilimitada e, portanto, divergente
(\cref{prop:b1:seq:first}): nenhuma \omterm{def:b1:seq:subsequence}{subsequência} convergente. Se $F$
não é \omterm{def:b1:topology:closed}{fechado}, tome $u_n \in F$ com $u_n \to \ell \notin F$
(\cref{thm:b1:topology:seqclosed}): toda \omterm{def:b1:seq:subsequence}{subsequência} converge para
$\ell \notin F$, de modo que \omterm{def:b1:seq:subsequence}{subsequência} alguma converge \emph{em} $F$.
\end{proof}

\begin{example}[Compactos encaixados]\label{ex:b1:topology:nested}
Um primeiro exercício para o teorema. Sejam $K_0 \supseteq K_1 \supseteq
K_2 \supseteq \dots$ subconjuntos compactos (\omterm{def:b1:topology:closed}{fechados} e limitados) não vazios
de $\R$. Então $\bigcap_n K_n \neq \emptyset$. Com efeito, tome $x_n
\in K_n$ para cada $n$: a sequência vive no compacto $K_0$,
de modo que uma \omterm{def:b1:seq:subsequence}{subsequência} $x_{\varphi(n)}$ converge para algum $x$
(\cref{thm:b1:topology:compact}). Para todo $m$ fixado, os termos
$x_{\varphi(n)}$ com $\varphi(n) \geq m$ estão todos no \omterm{def:b1:topology:closed}{fechado}
$K_m$, de sorte que o limite $x$ está em $K_m$
(\cref{thm:b1:topology:seqclosed}); e, como $m$ era arbitrário, $x \in
\bigcap_n K_n$. Um companheiro útil: se um \omterm{def:b1:topology:open}{aberto} $U$ contém
$\bigcap_n K_n$, então $U \supseteq K_n$ para algum $n$ --- aplique
o mesmo argumento a pontos $x_n \in K_n \setminus U$; o limite
$x$ estaria em $\bigcap K_n \subseteq U$, e, no entanto, $U$ \omterm{def:b1:topology:open}{aberto} força
$x_{\varphi(n)} \in U$ a partir de certo ponto, uma contradição. Os dois
enunciados falham sem compacidade: $\bigcap_n \intoo{0}{\frac
1n} = \emptyset$ e $\bigcap_n \intco{n}{+\infty} = \emptyset$.
A ideia de fechamento: a compacidade converte uma cadeia infinita de
afirmações de não vacuidade num único ponto-limite --- é a
ferramenta que sobrevive à passagem à interseção infinita, e o
problema de fim de semana (\cref{pb:b1:topology:1}) se apoiará nela duas vezes.
\end{example}

\begin{omfigure}
\begin{tikzpicture}[scale=9.5]
  \node[left, font=\small] at (-0.02,0) {$C_0$};
  \fill[omDef] (0,-0.011) rectangle (1,0.011);
  \node[left, font=\small] at (-0.02,-0.09) {$C_1$};
  \fill[omDef] (0,-0.101) rectangle (0.3333,-0.079);
  \fill[omDef] (0.6667,-0.101) rectangle (1,-0.079);
  \node[left, font=\small] at (-0.02,-0.18) {$C_2$};
  \fill[omDef] (0,-0.191) rectangle (0.1111,-0.169);
  \fill[omDef] (0.2222,-0.191) rectangle (0.3333,-0.169);
  \fill[omDef] (0.6667,-0.191) rectangle (0.7778,-0.169);
  \fill[omDef] (0.8889,-0.191) rectangle (1,-0.169);
  \node[left, font=\small] at (-0.02,-0.27) {$C_3$};
  \fill[omDef] (0,-0.281) rectangle (0.0370,-0.259);
  \fill[omDef] (0.0741,-0.281) rectangle (0.1111,-0.259);
  \fill[omDef] (0.2222,-0.281) rectangle (0.2593,-0.259);
  \fill[omDef] (0.2963,-0.281) rectangle (0.3333,-0.259);
  \fill[omDef] (0.6667,-0.281) rectangle (0.7037,-0.259);
  \fill[omDef] (0.7407,-0.281) rectangle (0.7778,-0.259);
  \fill[omDef] (0.8889,-0.281) rectangle (0.9259,-0.259);
  \fill[omDef] (0.9630,-0.281) rectangle (1,-0.259);
  \node[below, font=\small] at (0,-0.30) {$0$};
  \node[below, font=\small] at (0.3333,-0.30) {$\tfrac13$};
  \node[below, font=\small] at (0.6667,-0.30) {$\tfrac23$};
  \node[below, font=\small] at (1,-0.30) {$1$};
  \draw[omThm, thick, -Stealth] (0.5,-0.042) -- (0.5,-0.062);
  \draw[omThm, thick, -Stealth] (0.5,-0.132) -- (0.5,-0.152);
  \draw[omThm, thick, -Stealth] (0.5,-0.222) -- (0.5,-0.242);
\end{tikzpicture}

{\small Os quatro primeiros estágios da construção dos terços médios:
cada segmento de $C_n$ perde o seu terço médio \omterm{def:b1:topology:open}{aberto}, restando os
$2^{n+1}$ segmentos de $C_{n+1}$, de comprimento total
$(\frac23)^{n+1}$. O \omterm{pb:b1:topology:1}{conjunto de Cantor} $C = \bigcap_n C_n$ --- o
tema do problema de fim de semana \cref{pb:b1:topology:1} --- é
o resíduo compacto não vazio garantido pelo argumento dos compactos
encaixados acima: comprimento zero e, ainda assim, uma quantidade não
enumerável de pontos sobrevive.}
\end{omfigure}

\begin{remark}
Este é o motor por trás do teorema de Weierstrass
(\cref{ch:b1:continuity}) e do teorema de Heine sobre continuidade uniforme.
O nome ``compacto'' ganhará a sua definição geral (por coberturas)
no segundo ano; em $\R$, a compacidade sequencial é tudo de que
precisamos, e ``compacto $=$ \omterm{def:b1:topology:closed}{fechado} $+$ limitado'' é o enunciado a
guardar.
\end{remark}

\begin{example}[Fronteiras sob uniões]\label{ex:b1:topology:boundaryunion}
Sempre $\partial(A \cup B) \subseteq \partial A \cup \partial
B$: um ponto de $\partial(A \cup B)$ tem toda \omterm{def:b1:topology:open}{vizinhança}
encontrando $A \cup B$ (logo, $A$ ou $B$, um deles infinitas vezes)
e encontrando o complementar de $A \cup B$, que está contido nos
dois complementares --- uma verificação curta põe então o ponto em
$\partial A$ ou em $\partial B$. A inclusão pode ser espetacularmente
estrita: com $A = \Q$ e $B = \R\setminus\Q$,
\[
\partial(A \cup B) = \partial \R = \emptyset ,
\qquad
\partial A \cup \partial B = \R \cup \R = \R :
\]
dois \omterm{def:b1:logic:sets}{conjuntos} esfarrapados podem se colar num \omterm{def:b1:logic:sets}{conjunto} sem costuras,
com as suas \omterm{def:b1:topology:closure}{fronteiras} se aniquilando. A ideia de fechamento: os
\omterm{def:b1:topology:closure}{interiores} e os \omterm{def:b1:topology:closure}{fechos} se comportam \emph{monotonamente} sob uniões e
interseções, mas as \omterm{def:b1:topology:closure}{fronteiras} não --- trate $\partial$ como uma
grandeza derivada ($\overline A \setminus \mathring A$), nunca como um
operador com álgebra própria.
\end{example}

\begin{remark}[Perspectivas dentro deste volume]
O vocabulário aqui construído é consumido mais duas vezes neste livro.
No \cref{ch:b1:continuity}, todo teorema é um enunciado de topologia
disfarçado: o teorema do valor intermediário diz que as aplicações
contínuas preservam a propriedade de \omterm{prop:b1:reals:intervals}{intervalo}, e o teorema de
Weierstrass diz que elas preservam a compacidade --- e as demonstrações
chamam o \cref{thm:b1:topology:seqclosed,thm:b1:topology:compact} pelo nome. No
\cref{ch:b1:multivar}, as mesmas definições são relidas em $\R^2$, com
discos no lugar de \omterm{prop:b1:reals:intervals}{intervalos}: \omterm{def:b1:topology:open}{abertos}, \omterm{def:b1:topology:closure}{fechos} e compacidade se
transferem palavra por palavra, e o teorema de Weierstrass a duas
variáveis cavalga de novo sobre Bolzano--Weierstrass (extraia em cada
coordenada). A única noção que \emph{não} se generaliza sem dor é o
próprio \omterm{prop:b1:reals:intervals}{intervalo} --- no plano, a conexidade substitui a
convexidade, história que começa com o ``só $\emptyset$ e $\R$ são
\omterm{def:b1:topology:open}{abertos} e \omterm{def:b1:topology:closed}{fechados}'' do \cref{exo:b1:topology:9}.
\end{remark}

\section{Exercícios}

\begin{exercise}[$\star$]\label{exo:b1:topology:1}
Para cada \omterm{def:b1:logic:sets}{conjunto}, diga se ele é \omterm{def:b1:topology:open}{aberto}, \omterm{def:b1:topology:closed}{fechado}, os dois ou nenhum
(com justificativa): $\intoo{0}{1} \cup \intoo{2}{3}$; $\;\intco{0}{1}$;
$\;\{0\} \cup \intcc{1}{2}$; $\;\R \setminus \Z$; $\;\Q \cap
\intoo{0}{1}$.
\end{exercise}

\begin{exercise}[$\star$]\label{exo:b1:topology:2}
Determine $\mathring A$, $\overline A$ e $\partial A$ para:
$A = \intoc{0}{1} \cup \{2\}$; $\;A = \R \setminus \Q$;
$\;A = \bigl\{\frac{(-1)^n n}{n+1} : n \in \N\bigr\}$.
\end{exercise}

\begin{exercise}[$\star$]\label{exo:b1:topology:3}
Demonstre que um \omterm{def:b1:counting:card}{conjunto finito} é \omterm{def:b1:topology:closed}{fechado}, primeiro por complementos e
depois pela caracterização sequencial.
\end{exercise}

\begin{exercise}[$\star$]\label{exo:b1:topology:4}
Demonstre que, para quaisquer $A, B \subseteq \R$: $\overline{A \cup B} =
\overline A \cup \overline B$. Mostre com um exemplo que
$\overline{A \cap B}$ pode diferir de $\overline A \cap \overline
B$.
\end{exercise}

\begin{exercise}[$\star\star$]\label{exo:b1:topology:5}
Seja $u_n \to \ell$ em $\R$. Demonstre que o \omterm{def:b1:logic:sets}{conjunto} $\{u_n : n \in \N\}
\cup \{\ell\}$ é \omterm{def:b1:topology:closed}{fechado} (logo compacto, se acrescentarmos que ele é
limitado --- e ele é).
\end{exercise}

\begin{exercise}[$\star\star$]\label{exo:b1:topology:6}
Sejam $U$ \omterm{def:b1:topology:open}{aberto} e $A$ arbitrário. Demonstre que $U + A = \{u + a\}$ é
\omterm{def:b1:topology:open}{aberto}. Deduza que a soma de um \omterm{def:b1:topology:open}{aberto} com um \omterm{def:b1:logic:sets}{conjunto} qualquer é
\omterm{def:b1:topology:open}{aberta}, e contraste: exiba dois \omterm{def:b1:logic:sets}{conjuntos} \emph{\omterm{def:b1:topology:closed}{fechados}} cuja soma não
é \omterm{def:b1:topology:closed}{fechada}. \emph{(Tente $\Z$ e $\sqrt 2\,\Z$, com o \cref{exo:b1:reals:9}.)}
\end{exercise}

\begin{exercise}[$\star\star$]\label{exo:b1:topology:7}
Um ponto $x \in A$ é \emph{isolado} em $A$ quando alguma \omterm{def:b1:topology:open}{vizinhança} de
$x$ encontra $A$ apenas em $x$. Demonstre que todo ponto de $\Z$ é
isolado em $\Z$, que $A = \{\frac1n\}$ tem todos os seus pontos isolados e,
ainda assim, $\overline A \neq A$, e que um \omterm{def:b1:logic:sets}{conjunto} cujos pontos são todos
isolados tem \omterm{def:b1:topology:closure}{interior} vazio.
\end{exercise}

\begin{exercise}[$\star\star$]\label{exo:b1:topology:8}
Demonstre que o \omterm{def:b1:topology:closure}{fecho} de um \omterm{def:b1:logic:sets}{conjunto} limitado é limitado e que
$\sup \overline A = \sup A$ para $A$ não vazio e limitado superiormente. Deduza
que $\sup A \in \overline A$: o \omterm{def:b1:reals:bounds}{supremo} é sempre aderente.
\end{exercise}

\begin{exercise}[$\star\star\star$]\label{exo:b1:topology:9}
Demonstre que os únicos subconjuntos de $\R$ que são ao mesmo tempo
\omterm{def:b1:topology:open}{abertos} e \omterm{def:b1:topology:closed}{fechados} são $\emptyset$ e $\R$. \emph{Sugestão: suponha $A$
\omterm{def:b1:topology:open}{aberto}, \omterm{def:b1:topology:closed}{fechado}, com $A \neq \emptyset$ e $\R \setminus A \neq \emptyset$; tome $a \in
A$, $b \notin A$, digamos $a < b$, e considere $s = \sup\,(A \cap
\intcc{a}{b})$; decida se $s$ pode pertencer a $A$ ou ao seu
complementar.}
\end{exercise}

\begin{exercise}[$\star\star\star$]\label{exo:b1:topology:10}
(Estrutura dos \omterm{def:b1:topology:open}{abertos}) Seja $U \subseteq \R$ \omterm{def:b1:topology:open}{aberto} e não vazio. Para
$x \in U$, seja $I_x$ a união de todos os \omterm{prop:b1:reals:intervals}{intervalos} \omterm{def:b1:topology:open}{abertos} que contêm
$x$ e estão contidos em $U$. Demonstre que $I_x$ é um \omterm{prop:b1:reals:intervals}{intervalo} \omterm{def:b1:topology:open}{aberto},
que dois \omterm{def:b1:logic:sets}{conjuntos} $I_x$, $I_y$ são iguais ou disjuntos, e que $U$ é
uma união de \emph{enumeravelmente muitos} \omterm{prop:b1:reals:intervals}{intervalos} \omterm{def:b1:topology:open}{abertos} dois a
dois disjuntos \emph{(escolha um racional em cada um)}.
\end{exercise}

\begin{exercise}[$\star\star$]\label{exo:b1:topology:11}
Um ponto $x \in \R$ é \emph{ponto de acumulação} de $A$ quando
toda \omterm{def:b1:topology:open}{vizinhança} de $x$ encontra $A \setminus \{x\}$; o \omterm{def:b1:logic:sets}{conjunto} deles é
o \emph{\omterm{def:b1:logic:sets}{conjunto} derivado} $A'$. Demonstre que $\overline A = A \cup A'$
e que $A$ é \omterm{def:b1:topology:closed}{fechado} se, e somente se, $A' \subseteq A$. Determine
$A'$ para $A = \{\frac 1n : n \in \N^*\}$, para $A = \Z$ e para
$A = \Q$.
\end{exercise}

\begin{exercise}[$\star\star\star$]\label{exo:b1:topology:12}
Para $A \subseteq \R$ não vazio, defina $d_A(x) = \inf\{\abs{x - a}
: a \in A\}$. Demonstre:
\begin{enumerate}
  \item $\abs{d_A(x) - d_A(y)} \leq \abs{x - y}$ para todos $x, y$
        ($d_A$ é $1$-lipschitziana);
  \item $d_A(x) = 0$ se, e somente se, $x \in \overline A$; em
        particular, se $F$ é \omterm{def:b1:topology:closed}{fechado} e $x \notin F$, então
        $d_F(x) > 0$;
  \item para todo $\varepsilon > 0$, o \omterm{def:b1:logic:sets}{conjunto} $V_\varepsilon =
        \{x : d_F(x) < \varepsilon\}$ é \omterm{def:b1:topology:open}{aberto}, contém $F$ e
        $\bigcap_{\varepsilon > 0} V_\varepsilon = F$ para $F$
        \omterm{def:b1:topology:closed}{fechado}: todo \omterm{def:b1:topology:closed}{fechado} é uma interseção enumerável de
        \omterm{def:b1:topology:open}{abertos}.
\end{enumerate}
\end{exercise}

\section{Problema: O conjunto de Cantor, pequeno e enorme ao mesmo tempo}

\begin{problem}[{Problema de fim de semana --- o \omterm{pb:b1:topology:1}{conjunto de Cantor}
dos terços médios: comprimento zero, não enumerável, perfeito e $C + C =
\intcc{0}{2}$}]
\label{pb:b1:topology:1}
Retire de $\intcc{0}{1}$ o seu terço médio \omterm{def:b1:topology:open}{aberto}, depois o terço
médio de cada segmento restante, e repita para sempre: o que
sobrevive é o \emph{\omterm{pb:b1:topology:1}{conjunto de Cantor}} $C$, a fábrica fundamental de
contraexemplos da análise. Este problema o constrói, lê-o
pela maquinaria da base $3$ do \cref{pb:b1:reals:1}
e estabelece o seu retrato paradoxal: comprimento total zero e, ainda
assim, não enumerável; \omterm{def:b1:topology:closure}{interior} vazio e, ainda assim, sem ponto
isolado; totalmente desconexo e, ainda assim, $C + C$ preenche todo o
segmento $\intcc{0}{2}$.\index{conjunto de Cantor} Formalmente: $C_0 = \intcc{0}{1}$,
e $C_{n+1}$ se obtém de $C_n$ apagando o terço médio \omterm{def:b1:topology:open}{aberto} de cada
segmento de $C_n$; por fim, $C = \bigcap_{n \geq 0}
C_n$. Ao longo do problema, um
\emph{código ternário} de $x \in \intcc{0}{1}$ é qualquer cadeia de algarismos
$(d_k)_{k\geq1}$ com $d_k \in \{0, 1, 2\}$
cujo valor $\sup_n \sum_{k=1}^n d_k 3^{-k}$ seja igual a $x$ ---
os códigos impróprios (iguais a $2$ a partir de certo ponto) são
permitidos; pelo \cref{pb:b1:reals:1} (questões 9--11), todo $x \in
\intcc{0}{1}$ tem um ou dois códigos, dois exatamente quando $x = m/3^N
\in \intoo{0}{1}$.

\medskip
\noindent\textbf{Parte I --- A construção.}
\begin{enumerate}
  \item Descreva $C_1$ e $C_2$ explicitamente como uniões de
        segmentos e demonstre por indução: $C_n$ é uma união
        disjunta de $2^n$ segmentos \omterm{def:b1:topology:closed}{fechados}, cada um de comprimento
        $3^{-n}$.
  \item Mostre que $C$ é \omterm{def:b1:topology:closed}{fechado}, limitado --- logo compacto
        (\cref{thm:b1:topology:compact}) --- e não vazio, e que
        toda extremidade de todo segmento de todo $C_n$ pertence
        a $C$.
  \item O comprimento total de $C_n$ é $\bigl(\frac23\bigr)^n$.
        Deduza que, para todo $\varepsilon > 0$, o \omterm{def:b1:logic:sets}{conjunto} $C$ pode
        ser coberto por finitos segmentos de comprimento total
        $\leq \varepsilon$: o \omterm{pb:b1:topology:1}{conjunto de Cantor} tem \emph{comprimento
        zero}.
  \item Mostre que um \omterm{prop:b1:reals:intervals}{intervalo} contido em $C$ tem comprimento $\leq
        3^{-n}$ para todo $n$ e, portanto, é vazio ou um \omterm{def:b1:logic:sets}{conjunto}
        unitário: $\mathring C = \emptyset$. Sendo \omterm{def:b1:topology:closed}{fechado} de \omterm{def:b1:topology:closure}{interior}
        vazio, $C$ é \emph{nunca \omterm{def:b1:topology:dense}{denso}}.
\end{enumerate}

\medskip
\noindent\textbf{Parte II --- O código ternário.}
\begin{enumerate}
  \item Demonstre a recursão de autossemelhança
        \[
        C_{n+1} = \tfrac13 C_n \,\cup\,
        \bigl(\tfrac23 + \tfrac13 C_n\bigr),
        \qquad\text{logo}\qquad
        C = \tfrac13 C \,\cup\, \bigl(\tfrac23 + \tfrac13
        C\bigr),
        \]
        sendo as duas peças disjuntas: $C$ são duas cópias de si
        mesmo na escala $\frac13$.
  \item Demonstre por indução em $n$: $x \in C_n$ se, e somente se,
        $x$ tem um código ternário cujos $n$ primeiros algarismos estão
        em $\{0, 2\}$. Deduza, usando o fato de que $x$ tem no máximo
        dois códigos: $x \in C$ se, e somente se, $x$ tem um código
        \emph{sem nenhum algarismo igual a $1$} (um código
        \emph{sem $1$}).
  \item Códigos em ação: dê códigos sem $1$ para $0$, $1$,
        $\frac13$ e $\frac23$; mostre que $\frac14 = (0.\overline{02})_3$
        e $\frac34 = (0.\overline{20})_3$, de modo que ambos pertencem a
        $C$; e verifique que $\frac14$ \emph{não} é extremidade
        de nenhum $C_n$ (as extremidades têm a forma $m/3^n$).
  \item Mostre que cada $x \in C$ tem \emph{exatamente um} código
        sem $1$ \emph{(quando $x$ tem dois códigos, demonstre que
        exatamente um do par contém o algarismo $1$)}. Conclua: a
        \omterm{def:b1:logic:map}{aplicação} valor é uma bijeção das cadeias de $\{0,2\}$ sobre
        $C$.
  \item (Diagonal) Seja $k \mapsto x_k$ uma \omterm{def:b1:logic:map}{aplicação} qualquer $\N^* \to C$.
        Construa uma cadeia de $\{0, 2\}$ que difira no índice $k$ do
        código de $x_k$, e conclua que $C$ é não enumerável
        --- ao passo que, em contraste, a questão 3 diz que ele é
        metricamente desprezível.
\end{enumerate}

\medskip
\noindent\textbf{Parte III --- Retrato topológico.}
\begin{enumerate}
  \item Reúna o registro até aqui: $C$ é compacto, não enumerável,
        de comprimento zero e nunca \omterm{def:b1:topology:dense}{denso}. Que única inclusão
        $C \subseteq C_n$ carrega cada propriedade?
  \item ($C$ é perfeito) Seja $x \in C$ com código sem $1$
        $(d_k)$. Inverter o algarismo $d_n$ ($0 \leftrightarrow
        2$) produz $x_n \in C$ com $\abs{x_n - x} = 2 \cdot
        3^{-n}$. Conclua que $C$ não tem ponto isolado: todo
        ponto de $C$ é limite de \emph{outros} pontos de $C$.
  \item Mostre que as extremidades da questão 2 formam um subconjunto
        enumerável e \omterm{def:b1:topology:dense}{denso} de $C$ \emph{(trunque o código após $n$
        algarismos e continue com $0$s; enumerabilidade como no
        \cref{exo:b1:topology:10})}. Conclua: o ponto típico
        de $C$ --- como $\frac14$ --- \emph{não} é extremidade:
        as extremidades são um esqueleto enumerável dentro de um corpo
        não enumerável.
  \item (Totalmente desconexo) Sejam $x < y$ em $C$. Escolha $n$
        com $3^{-n} < y - x$ e produza um ponto $z \in
        \intoo{x}{y}$ com $z \notin C$. Conclua que os únicos
        \omterm{prop:b1:reals:intervals}{intervalos} não vazios contidos em $C$ são \omterm{def:b1:logic:sets}{conjuntos} unitários.
\end{enumerate}

\medskip
\noindent\textbf{Parte IV --- Aritmética de $C$.}
\begin{enumerate}
  \item Mostre que $1 - C = C$ \emph{(o que $x \mapsto 1 - x$ faz
        a um código sem $1$? lembre que $1 = (0.\overline{2})_3$)}.
  \item (Adição de códigos) Mostre que, se $x$ e $x'$ têm códigos
        $(a_k)$ e $(b_k)$, então $x + x' = \lim_n\,(t_n + t'_n)$,
        em que $t_n, t'_n$ são as somas parciais. Deduza: todo $y
        \in \intcc{0}{1}$ é o \emph{ponto médio} de dois pontos
        de $C$ --- dado um código $(e_k)$ de $y$, escolha algarismos
        $a_k, b_k \in \{0, 2\}$ com $\frac{a_k + b_k}{2} =
        e_k$.
  \item Conclua que $C + C = \intcc{0}{2}$ e, com a questão 14,
        que $C - C = \intcc{-1}{1}$. Instância concreta: escreva $1$ como
        soma dos dois pontos que não são extremidades encontrados na
        questão 7.
  \item Reflita: um \omterm{def:b1:logic:sets}{conjunto} de comprimento zero cujo \omterm{def:b1:logic:sets}{conjunto} das
        diferenças preenche $\intcc{-1}{1}$. Por que não há contradição
        entre ``$C$ é metricamente desprezível'' e ``$C + C$ tem
        comprimento pleno''? (Uma frase; pense no que o comprimento
        controla e no que ele não controla.)
\end{enumerate}

\medskip
\noindent\textbf{Parte V --- Membros, racionais e irracionais.}
\begin{enumerate}
  \item Combine a questão 8 com o critério de periodicidade do
        \cref{pb:b1:reals:1}: um ponto de $C$ é racional se, e somente
        se, o seu código sem $1$ é periódico a partir de certo ponto.
        Execute a divisão longa na base $3$ para verificar que $\frac1{13} =
        (0.\overline{002})_3 \in C$.
  \item Produza um membro explicitamente \emph{irracional} de $C$:
        o valor do código com $d_k = 2$ nas posições triangulares
        $k = \frac{j(j+1)}{2}$ e $d_k = 0$ nas demais. Justifique a
        irracionalidade pelo argumento dos espaçamentos crescentes do
        \cref{pb:b1:reals:1} (questão 20).
  \item (Sobre um segmento inteiro) Considere $h$ que leva o ponto de
        $C$ de código sem $1$ $(d_k)$ ao valor da cadeia
        \emph{binária} $\bigl(\frac{d_k}2\bigr)$, isto é,
        $h(x) = \sup_n \sum_{k=1}^n \frac{d_k}{2}\,2^{-k}$.
        Mostre que $h$ leva $C$ \emph{sobre} $\intcc{0}{1}$. Assim,
        o desprezível $C$ se projeta sobrejetivamente num segmento de
        comprimento pleno --- uma segunda demonstração de que $C$ é
        não enumerável.
  \item (Comprimento autossemelhante) Suponha que alguma noção de
        comprimento $L$ estivesse definida para $C$ e para as suas
        cópias encolhidas, respeitando a escala ($L(\lambda A) = \lambda L(A)$), a
        invariância por translação e a aditividade sobre a decomposição
        disjunta da questão 5. Mostre que então $L(C) = \frac23\,L(C)$,
        o que força $L(C) = 0$: só a autossemelhança já
        condena $C$ a comprimento zero.
\end{enumerate}

\medskip
\noindent\textbf{Parte VI --- Um \omterm{def:b1:arith:prime}{primo} gordo e a moral.}
\begin{enumerate}
  \item (\omterm{pb:b1:topology:1}{Conjunto de Cantor} gordo) Repita a construção, mas, no
        estágio $n$ ($n = 0, 1, 2, \dots$), retire de cada um dos $2^n$
        segmentos atuais apenas um \omterm{prop:b1:reals:intervals}{intervalo} \omterm{def:b1:topology:open}{aberto} central de
        comprimento $4^{-(n+1)}$. Mostre que os comprimentos $l_n$ dos
        segmentos obedecem a $l_{n+1} = \frac{l_n - 4^{-(n+1)}}{2}$, $l_n =
        \frac{2^n + 1}{2\cdot 4^n} > 0$, que o
        $K = \bigcap K_n$ resultante é compacto de \omterm{def:b1:topology:closure}{interior} vazio e
        que o comprimento total retirado é $\sum_{n\geq0} 2^n
        4^{-(n+1)} = \frac12$. Admitindo a aditividade (intuitiva, do terceiro ano) do
        comprimento para uniões finitas de \omterm{prop:b1:reals:intervals}{intervalos}, e
        usando os dois enunciados do \cref{ex:b1:topology:nested},
        mostre que toda família finita de \omterm{prop:b1:reals:intervals}{intervalos} \omterm{def:b1:topology:open}{abertos} que cubra
        $K$ tem comprimento total $\geq \frac12$: $K$ é nunca
        \omterm{def:b1:topology:dense}{denso}, mas \emph{não} é desprezível. Pequenez tem vários
        significados não equivalentes.
  \item (Distâncias) Mostre que, para um \omterm{def:b1:topology:closed}{fechado} não vazio $F \subseteq
        \R$ e $x \in \R$, o \omterm{def:b1:reals:bounds}{ínfimo} $d(x, F)$ é
        \emph{atingido} \emph{(sequência minimizante mais
        Bolzano--Weierstrass)}. Depois calcule
        \[
        \max_{y \in \intcc{0}{1}} d(y, C) = \frac16 ,
        \]
        atingido exatamente no centro $y = \frac12$ \emph{(um
        ponto de um vão criado no estágio $n$ está a menos de
        $\frac{3^{-n}}{2}$ das extremidades do vão, que estão em
        $C$)}.
  \item (Todo ponto é limite subsequencial) Usando as questões 12
        e 2, produza uma única sequência em $C$ cujo \omterm{def:b1:logic:sets}{conjunto} de
        limites subsequenciais seja \emph{todo} o $C$. (Compare:
        para uma sequência convergente esse \omterm{def:b1:logic:sets}{conjunto} é um único ponto
        --- $C$ realiza o extremo oposto entre os compactos.)
  \item Síntese, uma frase para cada: (i) que teoremas deste
        capítulo a construção de fato consumiu (estabilidade dos
        \omterm{def:b1:topology:closed}{fechados}, compacidade, caracterizações sequenciais)?
        (ii) liste os quatro pares paradoxais do retrato
        (comprimento zero/não enumerável, \omterm{def:b1:topology:closed}{fechado}/\omterm{def:b1:topology:closure}{interior} vazio,
        perfeito/totalmente desconexo, desprezível/$C+C$ pleno);
        (iii) onde $C$ ressurge mais adiante (a escada do diabo
        construída sobre $h$ na teoria da continuidade e a teoria da
        medida do volume do terceiro ano, em que $C$ separa
        ``enumerável'' de ``desprezível'')?
\end{enumerate}
\end{problem}
